\section{Experimental Setup}
\label{sec:setup}

\subsection{Testbed}

Most live results reported in Section~\ref{sec:evaluation} run on one
7-validator testbed: a genuine Cosmos SDK application on a forked CometBFT
core, deployed as a set of interconnected containers rather than an
in-process simulation. Two experiments are an explicit exception, carried
over from an earlier 4-validator cluster and not yet re-run at 7-validator
scale: the RQ1 malicious-proposer test (Section~\ref{sec:evaluation}) and
the eclipse/Sybil rows of Appendix~\ref{app:attacks}; we flag this
inconsistency rather than let the blanket ``7-validator'' framing elsewhere
in this section imply otherwise. The 7-validator testbed includes a
Bitcoin regtest network (two nodes, with a continuous mining loop producing
a block roughly every 20 seconds), a Celestia data-availability network
with seven independent bridge nodes -- one per validator, so each
validator retrieves its own DA view rather than the network sharing a
single bridge -- and a zero-knowledge prover (Noir, proved with
Barretenberg) that monitors for a \textsc{Sovereign}-mode backlog and
submits recovery-proof transactions once one accumulates. Validator gossip
runs over a dedicated full mesh of pairwise links, separate from the
network carrying Bitcoin, Celestia, and RPC traffic.

\subsection{Fault Injection}

Two classes of fault injection are used. Network-level faults are induced
through kernel-level traffic control, covering profiles for uniform
delay/jitter, packet loss, single-node isolation (an eclipse-style
partition), and heterogeneous per-validator wide-area-network delay/loss
intended to approximate seven geographically distinct regions.
Application-level faults are used where network delay alone cannot
reproduce the failure mechanism of interest: Bitcoin congestion is induced
by suspending new checkpoint submission at the application layer, so the
finality gap grows directly rather than through added RPC latency; DA
outages are induced by taking one or all seven Celestia bridge nodes
offline; and Byzantine behavior is induced either through dedicated
non-validator attacker processes or by configuring a validator instance to
misreport its proposals, double-sign, or flood the network with timeout
messages.

\subsection{Sim-to-Real Gap}
\label{subsec:sim-to-real}

This testbed approximates a real decentralized deployment; it is not one,
and we state the gap explicitly rather than leave it implicit. First, all
containers share a single physical machine's CPU and memory, so a
resource-contention condition can affect every validator simultaneously in
a way that geographically separate hardware would not; during evaluation
this manifested once as a contention-induced delay in the anchor
verification path, which we identified and corrected before the
measurements reported here were collected. Second, the testbed's Bitcoin
network produces blocks roughly 30$\times$ faster than Bitcoin mainnet;
height-based thresholds remain valid at any timescale, but the wall-clock
timings reported in Section~\ref{sec:evaluation} would be proportionally
longer in a mainnet deployment. Third, while the network runs seven
independent bridge nodes (one per validator, not a single shared bridge),
each validator's data-availability check currently trusts only its own
bridge rather than cross-validating against the other validators' bridges
through genuine light-client sampling; this is a deliberate scope choice
for the current prototype, and it means our DA-outage results measure a
bridge -- or all seven, taken offline together -- going down, not a
sampling-detected partial failure of an otherwise-healthy network. We
return to the implications of all three points in
Section~\ref{sec:discussion}.

\subsection{Baselines}

Three baselines are used across the experiments in
Section~\ref{sec:evaluation} and Appendix~\ref{app:attacks}: a vanilla
CometBFT deployment that never activates the FSM's health-aware hook and
therefore cannot leave its genesis \textsc{Anchored} state regardless of
real peripheral health, used to establish that the fallback behavior is
not an artifact of measurement; a peer-count-only P2P detector, used to
isolate the marginal benefit of the tri-interface profiler
(Section~\ref{subsec:sensors}); and each hysteresis edge's own floor
parameter value (a threshold of zero, disabling that edge's hysteresis),
used to isolate the marginal benefit of hysteresis itself.